\section{AI Agent 的部署与运营}

\subsection{成本与延迟 trade-off}

部署 Agent 不仅是模型选择，也是 resource orchestration：

\begin{itemize}
    \item \textbf{推理架构}：chat + tool loop vs. batch prompt，不同策略对 latency 影响巨大
    \item \textbf{资源阶梯}：热 path 使用 low-latency GPU，冷 path 使用 cheaper CPU
    \item \textbf{Auto-scaling}：按 request rate 自动扩缩容，避免 overprovisioning
    \item \textbf{Multi-model routing}：低风险任务走 smaller model，高风险任务走 Sonnet/Opus
\end{itemize}

\begin{importantbox}{不能只看吞吐，要看 end-to-end experience}
一个 Agent 的成功不是单纯推理速度，而是整体任务完成时间：LLM 回答→API 调用→后处理→用户确认。任何一步延迟都能破坏用户体验。Vertex AI 的 Agent Builder 通过 pipeline 仪表板，把每个组件 latency 具体归集到 dashboard。
\end{importantbox}

这里的经验是：Pod 规模和 prompt size 会构成两个相互挤压的维度。对于 token-heavy application，可以把 pipeline 分成「粗粒度 first pass」和「精细 second pass」，在保证 latency 的同时把 model call 费用最小化。

\subsection{数据与用户反馈闭环}

持续提升 Agent 需要循环利用 real user data：

\begin{itemize}
    \item \textbf{Structured feedback}：插入「是否解决问题」「是否合规」的按钮
    \item \textbf{Failure tagging}：把失败案例打上标签（ hallucination、tool failure、safety violation）
    \item \textbf{Retraining loop}：把高价值的 failure 推进 fine-tuning/LoRA
    \item \textbf{Prompt atlas}：每次 prompt 修改都版本化，便于回滚
\end{itemize}

\begin{knowledgebox}{Feedback loop 比「再训练」更早见效}
很多企业先用 user feedback 细调 prompt，再决定是否重新训练。这种「先微调交互，再大规模 retrain」策略减少了迭代成本，也能更快交付可信 Agent。
\end{knowledgebox}

在 Vertex AI 中，feedback loop 通常跑在独立的 `data labeling project` 中：将失败案例发给 labeler 标注后，自动触发 retrain pipeline。这种自动化闭环保证 model loop 时刻保持 fresh。

\subsection{本章小结}

Agent 运营侧重延迟成本、资源配比和数据闭环；把这些环节标准化后，团队才能把实验室能力推向生产环境。

